\documentclass[11pt]{article}

\usepackage[margin=1in]{geometry}
\usepackage[T1]{fontenc}
\usepackage{lmodern}
\usepackage{microtype}
\usepackage{amsmath,amssymb,mathtools}
\usepackage{booktabs}
\usepackage{array}
\usepackage{enumitem}
\usepackage{xcolor}
\usepackage{hyperref}
\usepackage[nameinlink,noabbrev]{cleveref}
\usepackage{listings}

\hypersetup{
  colorlinks=true,
  linkcolor=blue!55!black,
  citecolor=blue!55!black,
  urlcolor=blue!55!black,
  pdftitle={An explicit A6 point producer for Algorithm 5.1},
  pdfauthor={OpenAI}
}

\lstset{
  basicstyle=\ttfamily\small,
  columns=fullflexible,
  frame=single,
  breaklines=true,
  xleftmargin=0.5em,
  xrightmargin=0.5em,
  aboveskip=0.7em,
  belowskip=0.7em
}

\setlist[enumerate]{itemsep=0.35em,topsep=0.45em}
\setlist[itemize]{itemsep=0.25em,topsep=0.4em}
\setlength{\parindent}{0pt}
\setlength{\parskip}{0.55em}

\newcommand{\QQ}{\mathbf{Q}}
\newcommand{\ZZ}{\mathbf{Z}}
\newcommand{\PP}{\mathbf{P}}
\newcommand{\Aff}{\mathbf{A}}
\newcommand{\ol}[1]{\overline{#1}}
\newcommand{\Gal}{\operatorname{Gal}}
\newcommand{\Mat}{\operatorname{Mat}}
\newcommand{\Span}{\operatorname{Span}}
\newcommand{\im}{\operatorname{im}}
\newcommand{\rank}{\operatorname{rank}}
\newcommand{\Spec}{\operatorname{Spec}}
\newcommand{\cM}{\mathcal{M}}
\newcommand{\tcM}{\widetilde{\mathcal{M}}}
\newcommand{\tM}{\widetilde{M}}
\newcommand{\cO}{\mathcal{O}}
\newcommand{\cL}{\mathcal{L}}
\newcommand{\sI}{s}
\newcommand{\Gamhat}{\widehat{\Gamma}}
\newcommand{\gamhat}{\widehat{\gamma}}
\newcommand{\Deltav}{\Delta_{\mathrm{VdM}}}
\newcommand{\deltaE}{\delta_{E_6}}

\newenvironment{summarybox}{%
  \begin{center}\begin{minipage}{0.94\textwidth}\small
  \hrule\vspace{0.55em}}{%
  \vspace{0.55em}\hrule\end{minipage}\end{center}}

\title{An Explicit $\Aff^6$ Point Producer for the Twists\\
in Algorithm 5.1 of Elsenhans--Jahnel}
\author{Cuspidal-cubic parameters, Pinkham's $E_6$ action,\\
and integration with the descent data of Section 7}
\date{\today}

\begin{document}
\maketitle

\begin{abstract}
The rational-point search in step (vi) of Algorithm 5.1 can be replaced by an
explicit dominant rational map from affine six-space.  The construction puts six
ordered points on a fixed cuspidal cubic, uses Pinkham's linear reflection
representation of $W(E_6)$ on their six parameters, twists this six-dimensional
representation by the same Galois cocycle used in the paper, and evaluates Coble's
forty gamma coordinates.  A useful normalization by the six-variable Vandermonde
turns the restricted gamma coordinates into degree-nine polynomials and removes
the parameter-dependent scalar in the Cremona action.  This note gives all formulas,
the required linear systems, the conversion to the ten descended coordinates from
Section 7, and implementation-oriented pseudocode.  The resulting map is dominant,
so an exhaustive search in six integer parameters is guaranteed to produce a point
in the open twisted marked moduli space, not merely in its compactification.
\end{abstract}

\tableofcontents

\section{The conclusion and the formula to implement}

Retain the notation and the data from Section 7 of the earlier note.  Thus
$L/\QQ$ is finite Galois,
\[
  \rho:\Gal(L/\QQ)\xrightarrow{\sim}G\subseteq W(E_6),
\]
$J\in\Mat_{80\times10}(\QQ)$ identifies the ten-dimensional span of the
signed gamma coordinates with a subspace of $\QQ^{80}$, and
$B_\rho\in\operatorname{GL}_{10}(L)$ has columns forming a $\QQ$-basis of the
exact semilinear fixed space used to descend the ambient $\PP^9$.

There is a second fixed-space matrix
\[
  D_\rho\in\operatorname{GL}_6(L),
\]
computed from Pinkham's six-dimensional reflection representation.  It identifies
the twist of that representation with $\Aff^6_\QQ$:
\[
  u\in\QQ^6\longmapsto t=D_\rho u\in L^6.
\]
Let
\[
  \Gamhat:L^6\longrightarrow L^{80}
\]
be the explicitly evaluable signed Coble vector defined in
\cref{sec:normalized-gammas}.  It consists of degree-nine polynomials.  Choose a
rational left inverse
\[
  K\in\Mat_{10\times80}(\QQ),\qquad KJ=I_{10}.
\]
Then the desired map is
\begin{equation}
 \boxed{
 \Phi_\rho:\Aff^6_\QQ\dashrightarrow\tM_\rho\subseteq\PP^9_\QQ,
 \qquad
 u\longmapsto
 \left[B_\rho^{-1}K\Gamhat(D_\rho u)\right].
 }
 \label{eq:main-map}
\end{equation}
Although the vector inside brackets is written with entries in $L$, its coordinate
ratios lie in $\QQ$.  On the open set specified in \cref{sec:regular-locus}, the
image lies in the open twisted marked moduli space $\tcM_\rho$, hence represents a
smooth marked cubic surface.

\begin{summarybox}
\textbf{Algorithmic consequence.}
Steps (v) and (vi) of Algorithm 5.1---descending the thirty cubics and then searching
for a rational point on their common zero locus---are no longer needed for point
production.  One may retain the descended cubics as an optional certificate.  The
replacement consists only of a six-dimensional semilinear fixed-space computation,
evaluation of forty explicit degree-nine products, and linear algebra with $J$ and
$B_\rho$.
\end{summarybox}

\section{What must be added to the Section 7 data}

The construction reuses all of the ten-dimensional data already described in
Section 7.  Only three additions are needed.

\begin{enumerate}
\item A six-dimensional matrix representation
\[
  \mu:W(E_6)\longrightarrow\operatorname{GL}_6(\QQ),
  \qquad g\longmapsto M_g,
\]
with explicit generator matrices; see \cref{sec:pinkham-action}.

\item A basis matrix $D_\rho$ for the semilinear fixed space in $L^6$; see
\cref{sec:source-descent}.

\item An evaluator for the forty normalized gamma polynomials, in precisely the
same label order used to construct $J$ and the signed permutation matrices; see
\cref{sec:normalized-gammas,sec:ordering}.
\end{enumerate}

The safest universal data structure stores every element of $W(E_6)$ as a word in
fixed generators.  The same word is then evaluated in the $80$-dimensional signed
gamma representation, the ten-dimensional representation, and the six-dimensional
Pinkham representation.  This removes nearly all inverse and left-versus-right
action ambiguity.

\section{The untwisted cuspidal-cubic parameterization}
\label{sec:cusp-map}

\subsection{Six points on a fixed cuspidal cubic}

Let
\[
  C:\ yz^2=x^3\subseteq\PP^2_\QQ.
\]
Its smooth locus is parametrized by
\[
  p(t)=(t:t^3:1),\qquad t\in\Aff^1.
\]
For
\[
  t=(t_1,\ldots,t_6)\in\Aff^6,
\]
put
\[
  p_i(t)=p(t_i)=(t_i:t_i^3:1).
\]
Whenever these six ordered points are in general position, blowing them up and
using the anticanonical linear system produces a smooth cubic surface together with
the ordered six exceptional curves.  The latter are a marking.  Thus one obtains a
rational map
\[
  f:\Aff^6\dashrightarrow\tcM,
\]
where $\tcM$ denotes the moduli space of smooth marked cubic surfaces.

Duncan and Reichstein explain that this map is dominant: a general ordered
six-point configuration in $\PP^2$ lies on the smooth locus of some cuspidal cubic,
and all cuspidal plane cubics are projectively equivalent in characteristic different
from $2$ and $3$ \cite[Lemma 6.1]{DR-PR}.

\subsection{The regular locus}
\label{sec:regular-locus}

For a subset $I\subseteq\{1,\ldots,6\}$, write
\[
  s_I(t)=\sum_{i\in I}t_i,
  \qquad
  s_{[6]}(t)=t_1+\cdots+t_6.
\]
The six points are distinct precisely when
\[
  t_i-t_j\neq0\qquad(i\neq j).
\]
Three distinct points $p_i,p_j,p_k$ on $C$ are collinear precisely when
\[
  t_i+t_j+t_k=0,
\]
and all six lie on a conic precisely when
\[
  t_1+\cdots+t_6=0
\]
\cite[Lemma 6.1]{DR-PR}.  Consequently the general-position locus is
\[
  V^\circ=\{t\in\Aff^6:\deltaE(t)\neq0\},
\]
where
\begin{equation}
 \boxed{
  \deltaE(t)=
  \left(\prod_{1\leq i<j\leq6}(t_j-t_i)\right)
  \left(\prod_{1\leq i<j<k\leq6}(t_i+t_j+t_k)\right)
  (t_1+\cdots+t_6).
 }
 \label{eq:E6-discriminant}
\end{equation}
There are $15+20+1=36$ linear factors.  They are the positive-root hyperplanes for
the reflection representation of type $E_6$ used below.

On $V^\circ$, the map $f$ is a regular map to $\tcM$.  The complement of
$V^\circ$ is exactly the locus that must be rejected before evaluating a candidate
parameter vector.

\section{Pinkham's explicit six-dimensional $W(E_6)$ action}
\label{sec:pinkham-action}

\subsection{Simple reflections}

Regard $t=(t_1,\ldots,t_6)^t$ as a column vector.  Put
\[
  c=t_1+t_2+t_3.
\]
Pinkham's Cremona reflection is
\begin{equation}
\begin{aligned}
 t_i&\longmapsto t_i-\frac23c &&(i=1,2,3),\\
 t_j&\longmapsto t_j+\frac13c &&(j=4,5,6).
\end{aligned}
\label{eq:pinkham-reflection}
\end{equation}
In matrix form,
\begin{equation}
 M_{s_1}=\frac13
 \begin{pmatrix}
  1&-2&-2&0&0&0\\
 -2& 1&-2&0&0&0\\
 -2&-2& 1&0&0&0\\
  1& 1& 1&3&0&0\\
  1& 1& 1&0&3&0\\
  1& 1& 1&0&0&3
 \end{pmatrix}.
 \label{eq:pinkham-matrix}
\end{equation}
For $i=2,\ldots,6$, let $s_i$ interchange $t_{i-1}$ and $t_i$ and leave the
other coordinates fixed.  These six matrices satisfy the Coxeter relations of type
$E_6$.  With the numbering used here, the edges are
\[
  \{1,4\},\ \{2,3\},\ \{3,4\},\ \{4,5\},\ \{5,6\}.
\]
Thus
\[
 (s_is_j)^3=1
 \quad\text{for these five pairs},
 \qquad
 (s_is_j)^2=1
 \quad\text{for all other }i\neq j.
\]
Pinkham derives \eqref{eq:pinkham-reflection} from the standard quadratic
transformation centered at the first three blow-up points and proves that these
transformations give the Weyl reflection action \cite[pp.~196--198]{Pinkham}.
Duncan--Reichstein use this observation to make the dominant map
$\Aff^6\dashrightarrow\tcM$ $W(E_6)$-equivariant \cite[Lemma 6.1]{DR-PR}.

\subsection{The three generators matching the authors' Magma example}

The supplementary file \texttt{c9\_example.m} constructs $W(E_6)$ from three
operations: one Cremona transformation, a six-cycle of the points, and the
transposition of the first two points \cite{EJcode}.  The corresponding matrices on
$t$ are
\[
 M_{\mathrm{cr}}=M_{s_1},
\]
\[
 M_{\mathrm{cyc}}
 \begin{pmatrix}t_1\\t_2\\t_3\\t_4\\t_5\\t_6\end{pmatrix}
 =
 \begin{pmatrix}t_2\\t_3\\t_4\\t_5\\t_6\\t_1\end{pmatrix},
 \qquad
 M_{\mathrm{sw}}
 \begin{pmatrix}t_1\\t_2\\t_3\\t_4\\t_5\\t_6\end{pmatrix}
 =
 \begin{pmatrix}t_2\\t_1\\t_3\\t_4\\t_5\\t_6\end{pmatrix}.
\]
These three matrices generate the same six-dimensional representation.  If the
existing implementation represents elements of $G$ as words in the three gamma
matrices constructed from \texttt{bo -> bo2}, \texttt{bo -> bo3}, and
\texttt{bo -> bo4}, evaluate those identical words in
$M_{\mathrm{cr}},M_{\mathrm{cyc}},M_{\mathrm{sw}}$.

\subsection{Convention test}

The paper and its Magma code use row vectors in some places and column vectors in
others, and a permutation of coordinates can be represented either by a matrix or
its inverse.  Before any descent computation, perform the following universal test
for each chosen generator $g$ and a random $t\in V^\circ(\QQ)$:
\[
  [\Gamma(M_gt)]=[P_g\Gamma(t)]
  \quad\text{in }\PP^{79}.
\]
If instead the right side agrees with $P_g^{-1}\Gamma(t)$, replace every $M_g$ by
$M_{g^{-1}}$, or equivalently reverse the word convention.  The same adjustment
must then be used in the six- and ten-dimensional fixed equations.

\section{Coble coordinates restricted to the cusp}
\label{sec:normalized-gammas}

\subsection{Factorization of the basic determinants}

For $I\subseteq\{1,\ldots,6\}$, define
\[
  \Delta_I(t)=\prod_{\substack{i<j\\i,j\in I}}(t_j-t_i),
  \qquad
  \Deltav(t)=\Delta_{\{1,\ldots,6\}}(t).
\]
For a three-element set $I=\{i,j,k\}$, Coble's minor evaluated at
$p_r=(t_r:t_r^3:1)$ is
\begin{equation}
  m_I(t)=\Delta_I(t)s_I(t).
  \label{eq:minor-factor}
\end{equation}
Indeed, it is the determinant of the three rows
$(t_r,t_r^3,1)$ for $r\in I$.  The conic determinant of
Elsenhans--Jahnel, Definition 2.3, becomes
\begin{equation}
  d_2(t)=\Deltav(t)s_{[6]}(t).
  \label{eq:d2-factor}
\end{equation}
These identities turn every gamma coordinate into a product of linear forms.

\subsection{The ten coordinates of the first type}

For an unordered decomposition
\[
  I\sqcup I^c=\{1,\ldots,6\},\qquad |I|=|I^c|=3,
\]
Coble's first type is
\[
  \gamma_{I\mid I^c}=m_I m_{I^c}d_2.
\]
Using \eqref{eq:minor-factor} and \eqref{eq:d2-factor},
\[
  \gamma_{I\mid I^c}
  =\Deltav\,\Delta_I\Delta_{I^c}s_Is_{I^c}s_{[6]}.
\]
Define the normalized coordinate
\begin{equation}
 \boxed{
  \gamhat_{I\mid I^c}(t)
  :=\frac{\gamma_{I\mid I^c}(t)}{\Deltav(t)}
  =\Delta_I(t)\Delta_{I^c}(t)s_I(t)s_{I^c}(t)s_{[6]}(t).
 }
 \label{eq:gammahat-I}
\end{equation}
It is a homogeneous polynomial of degree $9$.

\subsection{The thirty coordinates of the second type}

Let $A,B,C$ be three unordered pairs partitioning $\{1,\ldots,6\}$.  A cyclic
orientation $(A,B,C)$ determines one of the two coordinates attached to the
partition.  Coble's formula is
\[
 \gamma_{A\to B\to C}
 =
 \prod_{a\in A}m_{\{a\}\cup B}
 \prod_{b\in B}m_{\{b\}\cup C}
 \prod_{c\in C}m_{\{c\}\cup A}.
\]
For a pair $A=\{a_1<a_2\}$, write
\[
  d_A=t_{a_2}-t_{a_1}.
\]
In the product above, every difference $t_j-t_i$ occurs at least once; the three
within-pair differences $d_A,d_B,d_C$ occur twice, and all twelve cross-pair
differences occur once.  Hence the difference part is
$\Deltav d_A d_B d_C$.  Therefore
\begin{equation}
\boxed{
\begin{aligned}
 \gamhat_{A\to B\to C}(t)
 &:=\frac{\gamma_{A\to B\to C}(t)}{\Deltav(t)}\\
 &=d_A d_B d_C
 \prod_{a\in A}s_{\{a\}\cup B}
 \prod_{b\in B}s_{\{b\}\cup C}
 \prod_{c\in C}s_{\{c\}\cup A}.
\end{aligned}}
\label{eq:gammahat-II}
\end{equation}
This is again a degree-nine polynomial.  Reversing the cyclic orientation from
$(A,B,C)$ to $(A,C,B)$ gives the other coordinate.

\subsection{The normalized gamma vector and the convention sign mask}

The formulas \eqref{eq:gammahat-I} and \eqref{eq:gammahat-II} use the
\emph{canonical minor convention}: whenever a minor $m_I$ occurs, its row indices
are first put in increasing order.  Choose exactly the same ordering of the forty
\emph{labels} as in the universal data and initially put
\[
  \Gamhat^{\mathrm{can}}_{40}(t)
  =\bigl(\gamhat_1(t),\ldots,\gamhat_{40}(t)\bigr)^t.
\]
Because every raw gamma coordinate has the common factor $\Deltav$ on the cusp,
\[
  [\Gamhat^{\mathrm{can}}_{40}(t)]=[\Gamma^{\mathrm{can}}_{40}(t)]
  \in\PP^{39}
  \qquad(t\in V^\circ).
\]
Thus the normalized vector gives exactly the same marked cubic surface.  Its lower
degree is computationally important: products of degree $9$ replace products of
degree $24$.

There is one small but essential implementation detail.  A stored implementation
may evaluate a determinant with its row indices in the order appearing in Coble's
product, rather than first sorting them.  This changes individual gamma coordinates
by fixed signs.  In particular, the authors' \texttt{gamma\_3} routine has this
feature for some of the thirty type-II coordinates.  Therefore store a universal
sign mask
\[
  \eta=(\eta_1,\ldots,\eta_{40})\in\{\pm1\}^{40}
\]
and define the normalized vector in the \emph{stored} forty-coordinate convention by
\begin{equation}
  \Gamhat^{\mathrm{st}}_{40}(t)
  =\bigl(\eta_1\gamhat_1(t),\ldots,
          \eta_{40}\gamhat_{40}(t)\bigr)^t.
  \label{eq:stored-sign-mask}
\end{equation}
The mask can be computed once and for all from a single regular rational test vector
$t^{(0)}$:
\begin{equation}
  \eta_a=
  \frac{\Gamma^{\mathrm{st}}_a(t^{(0)})}
       {\Deltav(t^{(0)})\gamhat_a(t^{(0)})}\in\{\pm1\}.
  \label{eq:sign-calibration}
\end{equation}
Equivalently, for each raw product multiply the parities required to sort the row
indices of its six minors.  The quotient in \eqref{eq:sign-calibration} must be
exactly $+1$ or $-1$ and should be asserted in code.  The ten type-I entries in the
authors' ordering have sign $+1$; among the thirty type-II entries some signs are
negative.

To pass to the paper's $80$ signed coordinates, use the stored signed-label list.
For each signed label $(\varepsilon,a)$, with $\varepsilon\in\{\pm1\}$ and
$1\leq a\leq40$, place
$\varepsilon\eta_a\gamhat_a(t)$ in the corresponding slot.  Denote the result by
\[
  \Gamhat(t)\in L^{80}.
\]
Do not assume that the universal order is ``all positive coordinates, then all
negative coordinates'' unless that is how $J$ and $P_g$ were constructed.  The
sign mask and the signed-label order are separate pieces of universal data.

\section{Why the Vandermonde normalization is especially useful}
\label{sec:exact-equivariance}

The map to projective gamma coordinates is already enough for descent.  The
normalization above has the additional advantage that it removes the nonconstant
common scalar in Pinkham's Cremona reflection.

\subsection{The root-factor substitution table}

Put
\[
  A_0=\{1,2,3\},\qquad B_0=\{4,5,6\},\qquad c=s_{A_0}.
\]
Under the reflection \eqref{eq:pinkham-reflection}:

\begin{center}
\renewcommand{\arraystretch}{1.18}
\begin{tabular}{>{\raggedright\arraybackslash}p{0.39\textwidth}
                >{\raggedright\arraybackslash}p{0.49\textwidth}}
\toprule
Linear form & Image \\
\midrule
$t_j-t_i$ with $i,j$ both in $A_0$ or both in $B_0$
  & unchanged \\
$t_j-t_i$ with $i\in A_0$, $j\in B_0$
  & $s_{(A_0\setminus\{i\})\cup\{j\}}$ \\
$s_{(A_0\setminus\{i\})\cup\{j\}}$ with $i\in A_0$, $j\in B_0$
  & $t_j-t_i$ \\
$s_{A_0}$ & $-s_{A_0}$ \\
$s_{B_0}$ & $s_{[6]}$ \\
$s_{[6]}$ & $s_{B_0}$ \\
A triple sum with one index in $A_0$ and two in $B_0$
  & unchanged \\
\bottomrule
\end{tabular}
\end{center}

Define
\[
 D_\times(t)=\prod_{i\in A_0}\prod_{j\in B_0}(t_j-t_i),
\]
\[
 N_\times(t)=
 \prod_{i\in A_0}\prod_{j\in B_0}
 s_{(A_0\setminus\{i\})\cup\{j\}},
 \qquad
 q(t)=\frac{N_\times(t)}{D_\times(t)}.
\]
The table gives
\begin{equation}
  \Deltav(M_{s_1}t)=q(t)\Deltav(t).
  \label{eq:vandermonde-transform}
\end{equation}
Substitution in Coble's forty formulas gives a signed permutation matrix $Q_1$ such
that
\[
  \Gamma_{40}(M_{s_1}t)=q(t)Q_1\Gamma_{40}(t).
\]
After division by $\Deltav$, the scalar cancels:
\begin{equation}
  \Gamhat_{40}(M_{s_1}t)=Q_1\Gamhat_{40}(t).
  \label{eq:exact-cremona}
\end{equation}

For an adjacent transposition $s_i$, $i=2,\ldots,6$, direct relabeling gives a
signed permutation $Q_i$ with
\[
  \Gamma_{40}(M_{s_i}t)=Q_i\Gamma_{40}(t),
  \qquad
  \Deltav(M_{s_i}t)=-\Deltav(t).
\]
Consequently
\begin{equation}
  \Gamhat_{40}(M_{s_i}t)=(-Q_i)\Gamhat_{40}(t).
  \label{eq:exact-symmetric}
\end{equation}
The six signed matrices
\[
  \widehat P_{s_1}=Q_1,
  \qquad
  \widehat P_{s_i}=-Q_i\quad(2\leq i\leq6)
\]
satisfy the $E_6$ Coxeter relations and yield an exact polynomial equivariance
\begin{equation}
  \Gamhat_{40}(M_gt)=\widehat P_g\Gamhat_{40}(t)
  \qquad(g\in W(E_6)).
  \label{eq:exact-equivariance}
\end{equation}

\subsection{How to construct the signed matrices without interpolation}

Equations \eqref{eq:gammahat-I} and \eqref{eq:gammahat-II} express each coordinate
as a multiset of linear root factors.  For each simple generator:

\begin{enumerate}
\item transform every factor using the table above or the corresponding adjacent
transposition;
\item sort the transformed factors;
\item locate the unique normalized gamma label with the same factor multiset;
\item record the sign coming from factors sent to their negatives.
\end{enumerate}

This produces the six $40\times40$ signed permutation matrices exactly.  It is a
symbolic alternative to the interpolation-and-matching step in the authors' Magma
example.  Doubling the signs in the usual way produces ordinary permutations on
the $80$ signed coordinates.

If $H=\operatorname{diag}(\eta_1,\ldots,\eta_{40})$ is the convention
sign mask, the matrices in the stored forty-coordinate convention are obtained by
conjugating by $H$.  The signed lift obtained this way and the signed lift stored in
an existing implementation can still differ by harmless global signs or by the
chosen ordering of the two signs, even though their projective actions agree.  For the map \eqref{eq:main-map}, only
projective agreement with the stored $P_g$ is required.  The generator test in
\cref{sec:pinkham-action} detects the necessary conversion.

\section{Twisting the six-dimensional source}
\label{sec:source-descent}

\subsection{The semilinear fixed space}

Let $g_1,\ldots,g_r$ be the same generators of $G$ used in the ten-dimensional
descent, and let
\[
  \sigma_i=\rho^{-1}(g_i)\in\Gal(L/\QQ).
\]
Suppose the convention in Section 7 imposes
\[
  R_{g_i}\sigma_i(y)=y
  \qquad\text{on }L^{10}.
\]
Use the identical convention on $L^6$:
\begin{equation}
  M_{g_i}\sigma_i(t)=t,
  \qquad i=1,\ldots,r.
  \label{eq:source-fixed}
\end{equation}
The common fixed vectors form a six-dimensional $\QQ$-vector space
\[
  V_\rho^{(6)}
  =\{t\in L^6:M_{g_i}\sigma_i(t)=t\text{ for every }i\}.
\]
This is the affine-space twist of Pinkham's reflection representation.

If the existing code instead uses the equivalent convention
$\sigma_i(y)=R_{g_i}y$, then use $\sigma_i(t)=M_{g_i}t$.  The point is not which
formula is preferred; the point is that the six- and ten-dimensional descents must
use the same descent datum.

\subsection{The rational linear system}

Choose a $\QQ$-basis $\omega_1,\ldots,\omega_n$ of $L$, and let
$S_{\sigma_i}\in\operatorname{GL}_n(\QQ)$ be the matrix of $\sigma_i$:
\[
  \sigma_i(\omega_b)=\sum_a(S_{\sigma_i})_{ab}\omega_a.
\]
Write
\[
  t_j=\sum_{a=1}^n U_{aj}\omega_a,
  \qquad U\in\Mat_{n\times6}(\QQ).
\]
Then \eqref{eq:source-fixed} is equivalent to
\begin{equation}
  S_{\sigma_i}U M_{g_i}^{t}=U.
  \label{eq:source-matrix-system}
\end{equation}
After vectorization,
\[
  (M_{g_i}\otimes S_{\sigma_i}-I_{6n})\operatorname{vec}(U)=0.
\]
Stack these equations for the selected generators.  The common kernel must have
$\QQ$-dimension $6$.

Convert a $\QQ$-basis of this kernel to six vectors
$d_1,\ldots,d_6\in L^6$ and put
\begin{equation}
  D_\rho=(d_1\ \cdots\ d_6)\in\operatorname{GL}_6(L).
  \label{eq:D-rho}
\end{equation}
Then
\[
  t=D_\rho u,\qquad u\in\QQ^6,
\]
parametrizes the twisted source.  Exact checks are
\[
  \det(D_\rho)\neq0,
  \qquad
  M_{g_i}\sigma_i(D_\rho)=D_\rho
  \quad\text{for every generator }g_i.
\]

\subsection{LLL reduction of the source lattice}

Coefficient growth can be reduced in the same way as in Remark 5.2 of the paper.
After computing $V_\rho^{(6)}$, form a full-rank lattice such as
\[
  \Lambda_\rho=V_\rho^{(6)}\cap\cO_L^6
\]
after clearing denominators and saturating.  Apply the Minkowski embedding and LLL
to obtain a short $\ZZ$-basis.  Use that basis as the columns of $D_\rho$.  This does
not alter the map; it merely replaces $u$ by a matrix in
$\operatorname{GL}_6(\QQ)$ and tends to make small integer vectors $u$ produce much
smaller number-field values.

\section{Converting a source point to the Section 7 coordinates}
\label{sec:coordinate-conversion}

\subsection{A left inverse for $J$}

Since $J$ has full column rank, choose
\[
  K\in\Mat_{10\times80}(\QQ),\qquad KJ=I_{10}.
\]
A sparse construction is preferable.  Choose ten row indices
$I=\{i_1,\ldots,i_{10}\}$ for which the submatrix $J_I$ is invertible.  Let
$E_I:\QQ^{80}\to\QQ^{10}$ select those entries.  Then
\[
  K=J_I^{-1}E_I.
\]
If the ten coordinates $y_1,\ldots,y_{10}$ were chosen to be ten of the signed gamma
coordinates, $J_I=I_{10}$ and $K$ is simply the coordinate selector.

\subsection{The conversion algorithm}

Given $u\in\QQ^6$, execute the following steps.

\begin{enumerate}[label=\textbf{P\arabic*.}]
\item Compute
\[
  t=D_\rho u\in L^6.
\]

\item Test the $36$ factors in \eqref{eq:E6-discriminant}.  If any factor vanishes,
reject $u$.

\item Evaluate the forty degree-nine polynomials
\eqref{eq:gammahat-I}--\eqref{eq:gammahat-II}, insert signs in the stored
$80$-coordinate order, and obtain
\[
  x=\Gamhat(t)\in L^{80}.
\]

\item Recover the original ten coordinates
\[
  y=Kx\in L^{10}.
\]
As an optional check, verify $Jy=x$.

\item Convert to the descended basis:
\[
  w=B_\rho^{-1}y\in L^{10}.
\]

\item Choose an index $j$ with $w_j\neq0$ and set
\begin{equation}
  z_i=\frac{w_i}{w_j},\qquad i=1,\ldots,10.
  \label{eq:rational-ratios}
\end{equation}
Every $z_i$ lies in $\QQ$.  Coerce them exactly to rational numbers and return
\[
  P=[z_1:\cdots:z_{10}]\in\tM_\rho(\QQ).
\]

\item For the subsequent Clebsch calculation, use the exact fixed representative
\[
  y_0=B_\rho z=\frac{y}{w_j},
  \qquad
  x_0=Jy_0=\frac{x}{w_j}.
\]
Extract the designated forty positive gamma entries from $x_0$.  Their power-sum
expressions give rational Clebsch invariants exactly as in the original algorithm.
\end{enumerate}

\subsection{Why the ratios are rational}

Let $\sigma\in\Gal(L/\QQ)$ and $g=\rho(\sigma)$.  Since
$t=D_\rho u$ is fixed by the twisted source action and the cusp map is
$W(E_6)$-equivariant, the gamma point satisfies the projective descent relation.
With the convention used here, there is a scalar $c_\sigma\in L^\times$ such that
\[
  P_g\sigma(x)=c_\sigma x.
\]
Because $x=Jy$ and $P_gJ=JR_g$, injectivity of $J$ gives
\[
  R_g\sigma(y)=c_\sigma y.
\]
Write $y=B_\rho w$.  The columns of $B_\rho$ are exact fixed vectors, so
\[
  R_g\sigma(B_\rho)=B_\rho.
\]
It follows that
\[
  \sigma(w)=c_\sigma w.
\]
Therefore
\[
  \sigma\left(\frac{w_i}{w_j}\right)=\frac{w_i}{w_j}
\]
for every $\sigma$, proving $w_i/w_j\in\QQ$.  This argument is why the raw gamma
vector need only satisfy projective, rather than vector, equivariance.

\section{The modified version of Algorithm 5.1}
\label{sec:modified-algorithm}

\subsection{Universal precomputation}

Compute once:

\begin{enumerate}
\item the forty gamma labels, the convention sign mask $\eta$, and the $80$
signed-label order;
\item $J$, the ten-dimensional matrices $R_g$, and optionally the thirty universal
cubic relations, as in the paper;
\item the six Pinkham generator matrices and a word map that evaluates the same
abstract $W(E_6)$ element in dimensions $80$, $10$, and $6$;
\item an evaluator for \eqref{eq:gammahat-I} and \eqref{eq:gammahat-II}
that applies the stored sign mask \eqref{eq:stored-sign-mask};
\item a left inverse $K$ of $J$.
\end{enumerate}

\subsection{Instance-specific computation}

Given $(L,G,\rho)$:

\begin{enumerate}[label=\textbf{Step \arabic*.}]
\item Compute the same generator/automorphism pairs $(g_i,\sigma_i)$ as in
Algorithm 5.1.

\item Compute $B_\rho$ from the ten-dimensional semilinear fixed equations, exactly
as in Section 7.

\item Evaluate the words for $g_i$ in the six-dimensional representation to obtain
$M_{g_i}$, and solve \eqref{eq:source-matrix-system} to obtain $D_\rho$.

\item Enumerate primitive vectors $u\in\ZZ^6$ by increasing height.  For each one,
apply steps P1--P6 of \cref{sec:coordinate-conversion}.  Reject $u$ if it lies on a
root hyperplane.

\item Optionally verify that the returned rational point $z$ satisfies the thirty
descended cubics.  These cubics are no longer needed to find $z$.

\item Use $x_0=JB_\rho z$ to compute the forty exact gamma values, the Clebsch
vector $[A:B:C:D:E]$, the pentahedral quintic, and the trace cubic of Algorithm
A.10.

\item If $E=0$, the pentahedral quintic is inseparable, or the resulting cubic is
singular, continue with the next vector $u$.
\end{enumerate}

\subsection{Language-neutral pseudocode}

\begin{lstlisting}
function PointFromA6(instance_data, universal_data):
    # Existing Section 7 objects:
    # B_rho, J, K, gamma_sign_mask, signed_gamma_order,
    # Clebsch formulas
    # New object:
    # D_rho, obtained from the 6-dimensional fixed equations

    for u in primitive_integer_vectors_by_height(6):
        t = D_rho * u                         # vector in L^6

        if any_E6_root_factor_is_zero(t):
            continue

        gamma40_can = canonical_normalized_cusp_gammas(t)
        gamma40 = apply_sign_mask(gamma40_can, gamma_sign_mask)
        x = insert_signed_coordinates(gamma40)
        y = K * x

        assert J * y == x

        w = inverse(B_rho) * y
        j = first_nonzero_coordinate(w)
        z = [w[i] / w[j] for i in 1..10]

        if not all_entries_are_exactly_rational(z):
            error("action convention or coordinate ordering mismatch")

        z = coerce_to_Q(z)
        x_exact = J * B_rho * z

        # Optional certificate:
        assert all(descended_cubic(z) == 0)

        clebsch = clebsch_from_positive_gammas(x_exact)
        if clebsch.E == 0:
            continue

        surface = algorithm_A10(clebsch)
        if surface_failed or surface_is_singular(surface):
            continue

        return z, surface
\end{lstlisting}

\section{Direct integration with \texttt{c9\_example.m}}
\label{sec:magma-integration}

The authors' example already contains nearly all universal objects:

\begin{itemize}
\item \texttt{all\_gamma(mat)} evaluates the forty raw gamma coordinates;
\item \texttt{bm} is a $10\times40$ row-basis matrix for their span;
\item \texttt{se} gives the forty linear forms in ten coordinates;
\item \texttt{mat\_l} contains the three $40\times40$ gamma matrices for the
Cremona generator, the six-cycle, and the transposition;
\item \texttt{z\_bm\_o\_10} and \texttt{para} encode the descended ten-dimensional
basis.
\end{itemize}

A minimally invasive patch has the following structure.

\begin{enumerate}
\item Store the three source matrices
$M_{\mathrm{cr}},M_{\mathrm{cyc}},M_{\mathrm{sw}}$ from
\cref{sec:pinkham-action}.

\item Whenever an element of the subgroup is constructed, retain a word in the
three generators.  Evaluate the word both in \texttt{mat\_l} and in the three source
matrices.  This gives the matching $40$- and $6$-dimensional matrices without a
group-identification problem.

\item Construct the six-dimensional semilinear fixed space in exact analogy with
the block-matrix construction used for the ten-dimensional lattice.  Its rank is six.
Apply LLL if desired and convert the result to $D_\rho$.

\item Delete the nested ten-variable point search.  Enumerate six integers, form
$t=D_\rho u$, and build the $6\times3$ matrix with rows
\[
  (t_i,t_i^3,1).
\]
The existing \texttt{all\_gamma} function may be used directly.  Divide its output
by
\[
  \prod_{i<j}(t_j-t_i)
\]
to obtain normalized values in exactly the authors' sign convention.  After this
version passes all checks, one may replace it by the degree-nine product formulas,
provided the universal sign mask \eqref{eq:stored-sign-mask} is applied.

\item Solve for the original ten coordinates using \texttt{bm}, then solve for the
descended coordinates using the matrix encoded by \texttt{para}.  Normalize one
nonzero descended coordinate to $1$ and verify that the remaining ratios are
rational.

\item Map the rational descended vector forward through \texttt{para} and
\texttt{se}.  Use this exact fixed gamma representative in the existing power-sum
and Clebsch code.
\end{enumerate}

\begin{summarybox}
\textbf{Row-vector warning.}
The paper's example uses expressions of the form
\texttt{row * bm} and therefore transposes several matrices relative to the column
convention in this note.  Translate the equations structurally, not by copying the
transpose symbols.  The two decisive tests are that the six-dimensional fixed space
has dimension $6$ and that the final descended ratios are fixed by every field
automorphism.
\end{summarybox}

\section{Existence, dominance, and termination of the search}

\subsection{Why every twist receives such a map}

The untwisted map
\[
  f:\Aff^6\dashrightarrow\tcM
\]
is dominant and $W(E_6)$-equivariant \cite[Lemma 6.1]{DR-PR}.  Twisting both sides
by the cocycle $\rho$ gives a dominant rational map
\[
  f_\rho:{}^\rho\Aff^6\dashrightarrow\tcM_\rho.
\]
The twist of a linear representation is a six-dimensional $\QQ$-vector space; the
matrix $D_\rho$ explicitly identifies it with $\Aff^6_\QQ$.  Formula
\eqref{eq:main-map} is the Coble-coordinate realization of this twisted map.
Dominance may be checked after base change to $L$, where it becomes the original
map composed with the invertible linear change of variables $D_\rho$.

It follows that
\[
  \tcM_\rho(\QQ)
\]
is Zariski dense.  In particular, the point sought in Algorithm 5.1 always exists.
The earlier bounded search could fail only because its chosen height bound was too
small.

\subsection{Why exhaustive integer enumeration terminates}

The pullback of the boundary is contained in the zero set of
\[
  \deltaE(D_\rho u)\in L[u_1,\ldots,u_6],
\]
a nonzero polynomial.  Hence rational vectors avoiding the boundary are Zariski
dense in $\Aff^6_\QQ$.

The conditions $E=0$, $\Delta=0$, and $F=0$ used in the paper define proper
divisors in the target.  Since $\Phi_\rho$ is dominant, their pullbacks are proper
closed subsets of $\Aff^6$.  Therefore an enumeration of primitive integer vectors
that is exhaustive by height must eventually find a vector avoiding the root
hyperplanes and all three downstream failure loci.  No useful small universal
height bound follows from this argument, but the search is no longer heuristic in
principle.

\section{Certification and debugging checklist}

A robust exact implementation should perform the following checks.

\subsection{Universal checks}

\begin{enumerate}
\item Verify the Coxeter relations for the six matrices $M_{s_i}$.
\item For random $t\in V^\circ(\QQ)$ and every universal generator, verify
\[
  [\Gamhat(M_gt)]=[P_g\Gamhat(t)].
\]
\item Verify the canonical degree-nine formulas against determinant evaluation:
\[
  \Gamma^{\mathrm{can}}_{40}(t)
  =\Deltav(t)\Gamhat^{\mathrm{can}}_{40}(t).
\]
Then verify the stored evaluator using the sign mask:
\[
  \Gamma^{\mathrm{st}}_{40}(t)
  =\Deltav(t)\Gamhat^{\mathrm{st}}_{40}(t).
\]
\item Verify $KJ=I_{10}$ and $P_gJ=JR_g$.
\item If constructing the normalized signed action symbolically, verify all Coxeter
relations for the resulting signed permutation matrices.
\end{enumerate}

\subsection{Instance checks}

\begin{enumerate}
\item The ten-dimensional fixed space has dimension $10$ and
$\det(B_\rho)\neq0$.
\item The six-dimensional fixed space has dimension $6$ and
$\det(D_\rho)\neq0$.
\item For every generator,
\[
  M_{g_i}\sigma_i(D_\rho)=D_\rho,
  \qquad
  R_{g_i}\sigma_i(B_\rho)=B_\rho,
\]
with the chosen convention.
\item For a candidate $u$, verify all $36$ root factors are nonzero.
\item Verify $Jy=x$ after $y=Kx$.
\item Verify all ratios in \eqref{eq:rational-ratios} lie in $\QQ$ exactly.  Failure
here almost always indicates an inverse, transpose, gamma-order, convention-sign-mask, or signed-label
mismatch.
\item Verify $x_0=JB_\rho z$ is exactly fixed under the $80$-coordinate descent.
\item If the thirty descended cubics were computed, verify that all vanish at $z$.
\item Verify the Clebsch quantities computed from $x_0$ are rational and that the
final cubic is smooth.
\item For a final independent check, compute the $27$ lines and their Galois action.
\end{enumerate}

\section{Practical refinements}

\subsection{Evaluate products, not determinants}

The formulas \eqref{eq:gammahat-I} and \eqref{eq:gammahat-II} need only additions,
subtractions, and products in $L$.  They avoid $6\times6$ determinants and reduce
the total degree from $24$ to $9$.  Precompute the following arrays for each $t$:
\[
  d_{ij}=t_j-t_i\quad(i<j),
\]
\[
  s_{ijk}=t_i+t_j+t_k\quad(i<j<k),
  \qquad s_{[6]}=t_1+\cdots+t_6.
\]
Every normalized gamma coordinate is then a product of nine entries from these
arrays.

\subsection{Search order}

After LLL reduction of the source lattice, enumerate primitive vectors by increasing
\[
  \|u\|_\infty=\max_i|u_i|
\]
or by increasing Euclidean norm.  Test the $36$ root factors before any gamma
evaluation.  Since the map is homogeneous of degree $9$ in $t$, scalar multiples of
$u$ give the same projective gamma point; primitive vectors avoid redundant work.

\subsection{Do not invert $B_\rho$ repeatedly}

Compute $B_\rho^{-1}$ once.  Better still, solve the linear system
$B_\rho w=y$ for each candidate using a stored decomposition.  The same applies to
$K$ if it was not chosen as a selector.

\subsection{Use the exact fixed representative downstream}

The raw source evaluation $x=\Gamhat(D_\rho u)$ is generally only projectively
fixed relative to the particular signed lift $P_g$ used in Section 7.  Do not coerce
its power sums individually to $\QQ$.  First compute $z$ from the coordinate ratios,
then replace $x$ by
\[
  x_0=JB_\rho z.
\]
This vector is exactly fixed and is the correct input to the Clebsch formulas.

\section{Canonical enumeration of the forty labels}
\label{sec:ordering}

A convenient abstract ordering is as follows.

\subsection{Ten type-I labels}

Enumerate all three-element subsets $I\subseteq\{1,\ldots,6\}$ containing $1$.
For each, use the unordered label $I\mid I^c$.  This gives ten labels.

\subsection{Thirty type-II labels}

Enumerate the fifteen perfect matchings of $\{1,\ldots,6\}$.  Sort the three pairs
lexicographically as $A<B<C$.  For each matching include the two cyclic
orientations
\[
  A\to B\to C,
  \qquad
  A\to C\to B.
\]
This gives thirty labels.

\subsection{The ordering in the authors' example}

The function \texttt{all\_gamma} in \texttt{c9\_example.m} places the thirty
type-II coordinates first and the ten type-I coordinates second.  In one-based
pseudocode, its type-II part is:

\begin{lstlisting}
for i,j,k,l,m in {2,...,6} with all five distinct:
    if j < k and l < m:
        A = {1,i}; B = {j,k}; C = {l,m}
        append gamma(A -> B -> C)
\end{lstlisting}

Its type-I part is:

\begin{lstlisting}
for i,j,k,l,m in {2,...,6} with all five distinct:
    if i < j and k < l and l < m:
        I  = {1,i,j}
        Ic = {k,l,m}
        append gamma(I | Ic)
\end{lstlisting}

These loops reproduce the authors' \emph{coordinate order}.  To reproduce their
coordinate signs as well, apply the mask \eqref{eq:stored-sign-mask}.  It can be
calibrated by \eqref{eq:sign-calibration}, or computed directly as follows: for each
entry returned by \texttt{gamma\_3}, take the product of the signs of the
permutations that sort the six row-index triples used by that routine.  Thus no
coordinate permutation is required, but the sign calibration must not be omitted.
The existing \texttt{all\_gamma} routine automatically supplies the correct signs;
dividing its output by $\Deltav$ is therefore the safest first implementation.

\section{A compact data-flow diagram}

The complete replacement for the nonlinear point search is
\[
\boxed{
\begin{array}{c}
 (g_i,\sigma_i),\ M_{g_i}
 \\[-1mm]
 \Downarrow\ \text{six-dimensional semilinear fixed equations}
 \\[-1mm]
 D_\rho\in\operatorname{GL}_6(L)
 \\[-1mm]
 \Downarrow\ u\in\ZZ^6\mapsto t=D_\rho u
 \\[-1mm]
 \Gamhat(t)\in L^{80}
 \\[-1mm]
 \Downarrow\ K
 \\[-1mm]
 y\in L^{10}
 \\[-1mm]
 \Downarrow\ B_\rho^{-1}\text{ and projective normalization}
 \\[-1mm]
 z\in\QQ^{10},\quad [z]\in\tcM_\rho(\QQ)
 \\[-1mm]
 \Downarrow\ JB_\rho z
 \\[-1mm]
 \text{exact gamma values}\to[A:B:C:D:E]\to\text{Algorithm A.10}.
\end{array}}
\]

\section{Final assessment}

The map can be made explicit enough for direct implementation.  The essential new
facts are:

\begin{enumerate}
\item Pinkham's reflection formula gives rational $6\times6$ matrices for the same
$W(E_6)$ elements already used in Algorithm 5.1.
\item On the cuspidal cubic, every Coble gamma coordinate has a common Vandermonde
factor.  Dividing it out gives the explicit degree-nine formulas
\eqref{eq:gammahat-I} and \eqref{eq:gammahat-II}.
\item Twisting the six-dimensional source is the same kind of rational linear
algebra already used for the ten-dimensional ambient space.
\item Formula \eqref{eq:main-map} converts any regular source parameter into a
rational point of the descended compactification; in fact it lands in the open
marked moduli space.
\item Dominance guarantees that exhaustive six-variable enumeration eventually
avoids every boundary and reconstruction failure divisor.
\end{enumerate}

The only delicate implementation issue is convention alignment among the abstract
Weyl-group word, its action on the six parameters, and its action on the signed
gamma coordinates.  The generator-level projective equivariance test and the
rational-ratio test isolate this issue immediately.

\begin{thebibliography}{9}

\bibitem{EJ}
A.-S. Elsenhans and J. Jahnel,
\emph{Moduli spaces and the inverse Galois problem for cubic surfaces},
Trans. Amer. Math. Soc. \textbf{367} (2015), no.~11, 7837--7861;
\href{https://arxiv.org/abs/1209.5591}{arXiv:1209.5591}.

\bibitem{DR-PR}
A. Duncan and Z. Reichstein,
\emph{Pseudo-reflection groups and essential dimension},
arXiv:1307.5724, especially Lemma~6.1.

\bibitem{DR-versality}
A. Duncan and Z. Reichstein,
\emph{Versality of algebraic group actions and rational points on twisted varieties},
J. Algebraic Geom. \textbf{24} (2015), 499--530;
\href{https://arxiv.org/abs/1109.6093}{arXiv:1109.6093}.

\bibitem{Pinkham}
H. Pinkham,
\emph{R\'esolution simultan\'ee de points doubles rationnels},
in \emph{S\'eminaire sur les singularit\'es des surfaces},
Lecture Notes in Math. \textbf{777}, Springer, 1980, 179--203;
see especially pp.~196--198.

\bibitem{EJcode}
J. Jahnel,
\emph{c9\_example.m: Magma example for Algorithm 5.1},
authors' supplementary webpage,
\url{https://wwwuser.gwdguser.de/~jjahnel/Arbeiten/Kub_Fl/c9_example.m}.

\end{thebibliography}

\end{document}
